% Independently audited and integrated as C.27.
\begin{theorem}[Proved: Kan assemblies over an exact base]
\label{C4:exact-base-kan-assembly}
Let $B$ be a nonempty connected small category exact over
$\mathbb Z$. Let $b\mapsto C_b=A(K_b)=\Delta_{\rm inj}/K_b$
be a diagram of categories and arbitrary functors, where every
$K_b$ is a nonempty connected small Kan complex. Put
$A=\int_B C_b$. The following conditions are equivalent:
\begin{enumerate}
\item $A$ is total;
\item $A$ is exact over $\mathbb Z$;
\item every positive homology diagram
$b\mapsto H_q(BC_b;\mathbb Z)$ is locally zero, and every
loop in every $BC_b$ becomes trivial after some base arrow.
\end{enumerate}
In the last condition, a base arrow acts by the specified
transition functor; the target basepoint is its image of the
given basepoint. In particular, allowing arbitrary nonfiltered
exact bases and arbitrary transition functors still produces
no total nonexact category by such an assembly.
\end{theorem}
\begin{proof}
Write $F_u:C_b\to C_c$ for the functor associated with
$u:b\to c$. The transition diagram is an ordinary strict
diagram; the Grothendieck construction has arrows
$(u,\alpha):(b,x)\to(c,y)$ with
$\alpha:F_u x\to y$.

\emph{The derived fibre formula, including nonconstant coefficients.}
For any diagram $N:A\to\Ab$ and the projection $\pi:A\to B$,
one has natural isomorphisms
\[
 (L_q\pi_!N)(b)=T_q^{C_b}(N|_{C_b}),\qquad
 T_q^A(N)=\operatorname{colim}_{b\in B}
             T_q^{C_b}(N|_{C_b}).
 \tag{EA.1}
\]
We spell out the first assertion to control coefficients.
An object of $\pi/b$ is $(c,x,u:c\to b)$. The fibre inclusion
$i:C_b\to\pi/b$ has left adjoint
\[
 R(c,x,u)=F_u x.
\]
The unit is the cocartesian arrow
$(c,x,u)\to(b,F_u x,\operatorname{id}_b)$, and $Ri=1$.
On bar complexes with the pulled-back coefficients, the map
induced by $R$ first applies $N$ to this unit at the initial
vertex of a bar string. Inclusion induces the reverse map.
One composite is the identity, and the other is chain homotopic
to the identity by the usual natural-transformation prism for
the unit. Thus these are coefficient-preserving derived
colimit comparisons, not only a constant-coefficient homology
calculation.

The pointwise derived left-Kan formula can be computed with
the same comma bar complex. For completeness, the restriction
to $\pi/b$ of a free representable $\mathbb Z[A(a_0,-)]$ is
colimit-acyclic: its category of elements splits according to
the composite arrow $\pi(a_0)\to b$, and each part has the
initial object given by $a_0$, that base arrow, and
$\operatorname{id}_{a_0}$. Projective resolutions therefore
justify the pointwise formula for arbitrary $N$. Finally
$\operatorname{colim}_A=\operatorname{colim}_B\pi_!$.
Since the second colimit is exact, it commutes with homology
of a projective resolution, giving the second equality in
(EA.1).

\emph{Persistent homology contradicts totality.}
Every category of elements of a set action on $B$ is exact:
its left Kan extension to $B$ is the exact fibrewise direct
sum, as checked in Theorem~\ref{C4:grothendieck-exclusion}.
Let $H(b)=H_q(BC_b;\mathbb Z)$ for $q>0$. If an element is
never killed by a base arrow, move it, if necessary, to have
the least finite order attainable by an image; if no finite
order is attainable, all its images have infinite order.
Thus for some $n\in\{2,3,\ldots,\infty\}$ the set action
\[
 T(b)=\{v\in H(b):\operatorname{ord}(H(u)v)=n
                    \text{ for every }u\text{ out of }b\}
\]
is nonempty somewhere. On a nonempty connected component
$D$ of $\operatorname{El}_B T$, the pulled-back homology
diagram contains the constant cyclic group $C_n$ injectively,
via $1\mapsto v$; here $C_\infty=\mathbb Z$. Exactness of
$D$ gives a nonzero injection of $C_n$ into its colimit.
Pull $T$ back along $A\to B$. The action component over $D$
is $\int_D C_b$, and is connected because the base and fibres
are connected. Formula (EA.1), now over $D$ and with constant
coefficients, shows that its $q$th homology is
$\operatorname{colim}_D H(b)\ne0$. Totality is contradicted.
This proves the homological part of (iii).

\emph{Persistent fundamental-group loops contradict totality.}
Let $\Pi_b$ be the fundamental groupoid of $BC_b$, on the
vertices that are objects of $C_b$. The functors $F_u$ induce
strict functors of these groupoids. Suppose that a loop
$g\in\Pi_b(z,z)$ is never sent to the identity by a base arrow.
As in the preceding paragraph, move to a loop whose order
$n$ persists under every further base arrow. Define
\[
 T(b)=\{(z,g): z\in C_b,\ g\in\Pi_b(z,z),
       \ \operatorname{ord}(F_u g)=n
          \text{ for every }u\text{ out of }b\}.
 \tag{EA.2}
\]
This is a nonempty covariant set action on $B$.
There is a covariant set action $L$ on the whole assembly,
defined at $(b,x)$ by
\[
 L(b,x)=\coprod_{(z,g)\in T(b)}
      \Pi_b(z,x)/\langle g\rangle.
 \tag{EA.3}
\]
The quotient in (EA.3) is by right composition: $h$ and
$h\circ g^k$ represent the same element. For
$(u,\alpha):(b,x)\to(c,y)$ the action sends
\[
 (z,g,[h])\longmapsto
       (F_u z,F_u g,[\alpha\circ F_u h]).
 \tag{EA.4}
\]
Right powers map to right powers, so (EA.4) is well defined;
strict composition of the transition functors and left
composition by fibre arrows verify the action law.

For a fixed $(z,g)$ the corresponding component of the fibre
action category is the connected cover of $BC_b$ with
fundamental group $\langle g\rangle$. Indeed the fibre
action is the usual path-groupoid coset action, based at
$(z,[\operatorname{id}_z])$; all fibre arrows act bijectively.
Its first homology is canonically $C_n$, with $1$ represented
by $g$. A base arrow carries this chosen basepoint and loop
to the chosen basepoint and loop for $(F_u z,F_u g)$.
Since their orders remain $n$, the map on these first
homology groups is the identity under the stated identifications.

Choose a connected component $D$ of $\operatorname{El}_B T$.
The corresponding part of $\operatorname{El}_A L$ is a
Grothendieck construction over the exact category $D$, with
these connected covering categories as fibres. It is connected.
Formula (EA.1) gives its first homology as
$\operatorname{colim}_D C_n=C_n\ne0$. This contradicts
totality. Therefore every fibre loop is killed by some base
arrow, completing (i)$\Rightarrow$(iii).

\emph{A sequence that kills a given derived-colimit cycle.}
Assume (iii), and fix $N:A\to\Ab$, a base object $b_0$, and
a positive-degree cycle in the bar complex of
$(C_{b_0},N|_{C_{b_0}})$. We construct a sequence of base
arrows and countable connected Kan subcomplexes
\[
 b_0\longrightarrow b_1\longrightarrow b_2\longrightarrow\cdots,
 \qquad K'_j\subset K_{b_j},
 \qquad C'_j=A(K'_j),
 \tag{EA.5}
\]
such that every transition functor restricts to
$C'_j\to C'_{j+1}$, the initial cycle lies in $C'_0$, and
the categorical colimit $C_\infty=\colim_j C'_j$ has weakly
contractible classifying space.

Start with all labels occurring in the finitely many objects
of the given cycle. They lie in a countable connected Kan
subcomplex: close under faces and degeneracies, add paths to
one vertex, choose horn fillers, and repeat. Each $C'_j$ is
countable. Enumerate its positive integral bar cycles and
its fundamental-groupoid loops. At step $j$, schedule the
first $j$ entries in the enumerations for the first $j$
stages. Their images at the current stage are finitely many
homology classes and loops. By (iii), each is killed by a
base arrow. Exactness of $B$ supplies condition I, so finitely
many such arrows can be amalgamated to one common composite
$b_j\to b_{j+1}$ that kills all scheduled classes and loops.

In the full category $C_{b_{j+1}}$, choose finite bar-boundary
witnesses for the killed cycles and finite nerve nullhomotopies
for the killed loops. These use only finitely many object
labels for each scheduled entry. Choose $K'_{j+1}$ containing
those labels and the labels of the images of every object of
$C'_j$. This is still countable; close it under faces,
degeneracies, connecting paths, and horn fillers as before.
The transition functor now restricts as required, and all
scheduled witnesses lie in $C'_{j+1}$. This construction works
even when the full $K_b$ are uncountable and the transition
functors are not induced by simplicial maps.

Every positive homology class and every loop from any stage
eventually occurs in the schedule and becomes zero. Filtered
categorical colimits commute with nerves and with the finite
homotopy and homology tests used here, as checked in
Theorem~\ref{C:kan-filtered-totalization}. Hence $BC_\infty$
is connected, simply connected, and integrally acyclic, so it
is weakly contractible. That theorem now implies that
$C_\infty$ is integrally exact.

Restrict $N$ to the categories in (EA.5). The cocartesian arrows
of $A$ supply compatible natural maps from each coefficient
diagram to the pullback of the next. They define a coefficient
diagram $N_\infty$ on $C_\infty$: for an object represented
at stage $j$, take the filtered colimit of its groups along
the succeeding images. Two representatives of the same object
agree at a later stage, so the definition is independent of
the representative, and the same argument defines its arrow
maps. Degree by degree,
\[
 \operatorname{colim}_j
    C_*(C'_j;N|_{C'_j})
       \cong C_*(C_\infty;N_\infty).
 \tag{EA.6}
\]
Indeed each bar string, each coefficient, and each equality
involves finitely many data and is represented at one stage.
The original positive-degree cycle is a boundary on the right
of (EA.6), by exactness of $C_\infty$. It is therefore a
boundary at one finite stage on the left. Thus, for every
$q>0$, the diagram $b\mapsto T_q^{C_b}(N|_{C_b})$ is locally
zero.

Formula (EA.1) now gives $T_q^A(N)=0$ for $q>0$, because
the colimit of a locally zero diagram is zero. This proves
(iii)$\Rightarrow$(ii). Integral exactness implies totality,
so (ii)$\Rightarrow$(i) completes the proof.
\end{proof}
